\documentclass[10pt,twocolumn]{article}
\usepackage[UTF8]{ctex}          % 中文支持
\usepackage{amsmath,amssymb}
\usepackage{booktabs}
\usepackage{multirow}
\usepackage{graphicx}
\usepackage{hyperref}
\usepackage{geometry}
\geometry{a4paper, margin=1in}

% 如在 Overleaf 编译，选择 XeLaTeX 引擎

\title{SimVLA：面向机器人操作的简洁视觉-语言-动作基线\\
\large 算法与实验（草稿）}
\author{}
\date{}

\begin{document}
\maketitle

% ============================================================
\section{方法}
% ============================================================

\subsection{整体架构}

SimVLA 是一个面向机器人操作任务的视觉-语言-动作（Vision-Language-Action, VLA）模型，
由三个核心模块构成：视觉语言主干网络、动作 Transformer 头以及基于流匹配的训练范式。

\noindent\textbf{视觉语言主干。}
我们采用 SmolVLM-500M-Instruct 作为多模态感知模块。
给定 $V$ 个相机视角的图像观测 $\{I_v\}_{v=1}^{V}$（$V{=}3$，包含正面视角、俯视视角和腕部视角）
以及自然语言任务指令 $\ell$，主干网络首先通过 SigLIP 视觉编码器提取图像特征，
再经多模态投影层映射至语言模型的特征空间，最终由语言模型对拼接后的视觉-语言
token 序列进行前向传播，得到融合表示
$\mathbf{F} \in \mathbb{R}^{B \times T_{\mathrm{enc}} \times D_{\mathrm{vlm}}}$。

\noindent\textbf{本体感知编码。}
机器人末端执行器的本体感知状态 $s \in \mathbb{R}^{8}$
（末端位置3维、末端姿态3维、夹爪状态2维）经 z-score 归一化后，
作为辅助条件输入动作 Transformer。

\noindent\textbf{动作 Transformer 头。}
动作头以 $\mathbf{F}$、本体感知 $s$ 和噪声化动作序列 $x_t$ 为输入，
预测流速度场，并通过 Euler 积分在推理时生成动作轨迹（详见第~\ref{sec:transformer} 节）。

% ------------------------------------------------------------
\subsection{动作 Transformer}
\label{sec:transformer}

动作 Transformer $f_\theta$ 采用标准 Pre-LayerNorm Transformer 结构，
负责在流匹配框架下预测速度场。其输入序列由以下 token 拼接而成：
\begin{equation}
  x = \mathrm{Concat}\!\bigl[
    \mathrm{ActionEnc}(x_t \oplus s \oplus \tau_t),\;
    \mathbf{W}_{\mathrm{vlm}}\mathbf{F}
  \bigr] + \mathbf{E}_{\mathrm{pos}}
\end{equation}
其中 $x_t \in \mathbb{R}^{B \times T_a \times 7}$ 为当前流时刻的噪声化动作序列，
$T_a$ 为动作预测时域长度，$\tau_t$ 为 sinusoidal 时间嵌入，
$\mathbf{W}_{\mathrm{vlm}}$ 为将 VLM 特征投影至动作 Transformer 隐层空间的线性层，
$\mathbf{E}_{\mathrm{pos}}$ 为可学习位置编码。
经 $L$ 层 Transformer block 后，提取前 $T_a$ 个位置的输出特征
$\mathbf{H} \in \mathbb{R}^{B \times T_a \times d}$，
再经解码器映射为速度预测 $\hat{v}_t$。

% ------------------------------------------------------------
\subsection{流匹配训练}
\label{sec:flow}

我们采用条件流匹配（Conditional Flow Matching）框架对动作分布进行建模。
给定归一化的目标动作 $a_0 \in \mathbb{R}^{T_a \times 7}$ 和
标准高斯噪声 $\epsilon \sim \mathcal{N}(0, I)$，
定义流时刻 $t \in (0,1)$ 处的插值轨迹为：
\begin{equation}
  x_t = t \cdot \epsilon + (1 - t) \cdot a_0
\end{equation}
对应的目标速度场为 $u_t = \epsilon - a_0$。
模型以 $x_t$、$t$、本体感知 $s$ 和 VLM 特征 $\mathbf{F}$ 为条件，
预测速度 $\hat{v}_t = f_\theta(x_t, t, s, \mathbf{F})$，训练目标为：
\begin{equation}
  \mathcal{L} = \mathbb{E}_{t,\,a_0,\,\epsilon}
    \bigl[\|\hat{v}_t - u_t\|_2^2\bigr]
\end{equation}
时间 $t$ 从偏向高噪声端的 Beta 分布中采样：
$t \sim \mathrm{Beta}(1.5,\, 1.0) \times 0.999 + 0.001$。

\noindent\textbf{推理。}
从高斯噪声 $x_1 \sim \mathcal{N}(0, I)$ 出发，
执行 $N_{\mathrm{step}}$ 步 Euler 积分：
\begin{equation}
  x_{t - \Delta t} = x_t + \Delta t \cdot f_\theta(x_t, t, s, \mathbf{F}),
  \quad \Delta t = -\tfrac{1}{N_{\mathrm{step}}}
\end{equation}
最终 $x_0$ 经反归一化后即为预测动作序列。

% ------------------------------------------------------------
\subsection{CTAF：连续时间动作场}
\label{sec:ctaf}

\noindent\textbf{动机。}
标准动作解码器对每个时间步独立预测一个动作向量，
导致预测轨迹缺乏时间维度的全局一致性，相邻帧之间可能产生不平滑的跳变，
不利于实际机器人执行。

\noindent\textbf{方法。}
我们提出连续时间动作场（Continuous-Time Action Field, CTAF），
将离散的 per-token 解码替换为傅里叶系数解码，使输出轨迹具有 $C^\infty$ 光滑性。

具体地，将 $T_a$ 个动作特征沿时间维度做均值池化，
得到全局特征向量 $\bar{\mathbf{h}} \in \mathbb{R}^d$，
再通过系数解码器映射为傅里叶系数张量：
\begin{equation}
  \mathbf{C} = \mathbf{W}_{\mathrm{coeff}}\,\bar{\mathbf{h}}
  \;\in\; \mathbb{R}^{(2M-1) \times D_a}
\end{equation}
其中 $M$ 为频率数（默认 $M{=}5$），$D_a{=}7$ 为动作维度。
系数布局为：第0行为直流分量 $c_0$，
第 $2k{-}1$ 行和第 $2k$ 行分别为频率 $k$ 的余弦和正弦系数
（$k=1,\ldots,M{-}1$）。

在 $T_a$ 个均匀时间点 $\tau_i = i/(T_a{-}1) \in [0,1]$ 处重建轨迹：
\begin{equation}
  a(\tau) = c_0 + \sum_{k=1}^{M-1}
    \bigl[c_k^{\cos}\cos(2\pi k\tau)
        + c_k^{\sin}\sin(2\pi k\tau)\bigr]
\end{equation}
以矩阵形式写为
$\hat{v}_t = \mathbf{B}\,\mathbf{C} \in \mathbb{R}^{T_a \times D_a}$，
其中 $\mathbf{B} \in \mathbb{R}^{T_a \times (2M-1)}$ 为傅里叶基矩阵（推理前预计算）。

\noindent\textbf{特性。}
CTAF 用 $\mathcal{O}(MD_a)$ 个参数描述一条完整轨迹，
轨迹光滑性由频率截断隐式保证，无需额外正则化。

% ------------------------------------------------------------
\subsection{PSCA：物理自洽适配}
\label{sec:psca}

\noindent\textbf{动机。}
大规模预训练 VLM 骨干与轻量级动作头之间存在特征分布差距，
全量微调成本高昂且易引发灾难性遗忘。
我们希望以极小的参数代价增强动作头对物理约束的适应能力。

\noindent\textbf{方法。}
我们提出物理自洽适配（Physical Self-Consistency Adaptation, PSCA），
在动作 Transformer 的每个 MLP 块中引入低秩 LoRA 适配器。

设标准两层 MLP 为
$\mathrm{fc}_1{:}\ \mathbb{R}^{d} \to \mathbb{R}^{4d}$，
$\mathrm{fc}_2{:}\ \mathbb{R}^{4d} \to \mathbb{R}^{d}$，
PSCA 在每层叠加 LoRA 残差：
\begin{align}
  h   &= \mathrm{GELU}\!\bigl(\mathbf{W}_1 x + \mathbf{B}_1\mathbf{A}_1 x\bigr) \\
  \mathrm{out} &= \mathbf{W}_2 h + \mathbf{B}_2\mathbf{A}_2 h
\end{align}
其中 $\mathbf{A}_1 \in \mathbb{R}^{r \times d}$，
$\mathbf{B}_1 \in \mathbb{R}^{4d \times r}$，
$r$ 为 LoRA 秩（默认 $r{=}8$）。
初始化时令 $\mathbf{B}_1 = \mathbf{B}_2 = \mathbf{0}$，
确保训练起始阶段 $\Delta W = \mathbf{B}\mathbf{A} = \mathbf{0}$，
不破坏已收敛的基础权重。

\noindent\textbf{测试时适配。}
推理阶段 PSCA 支持仅更新 $\mathbf{A},\mathbf{B}$ 参数而冻结基础权重，
利用物理一致性误差信号进行快速在线适配，无需重新训练整个模型。

\noindent\textbf{参数开销。}
设动作 Transformer 有 $L$ 层，每层 MLP 引入的额外参数量为
$2r(d + 4d) = 10rd$。
对于大模型（$d{=}1024,\ L{=}24,\ r{=}8$），
PSCA 额外引入约 1.97M 参数，相比整体参数量（$\approx$350M）占比不足 0.6\%。

% ------------------------------------------------------------
\subsection{模型配置}
\label{sec:config}

\begin{table}[h]
\centering
\caption{SimVLA 模型配置}
\begin{tabular}{lcccc}
\toprule
配置 & 隐层维度 & 层数 & 头数 & 参数量（动作头）\\
\midrule
SimVLA-Small & 768  & 12 & 12 & $\approx$85M \\
SimVLA-Large & 1024 & 24 & 16 & $\approx$350M \\
\bottomrule
\end{tabular}
\end{table}

所有配置共享 SmolVLM-500M-Instruct 骨干（$\approx$500M 参数），
输入分辨率为 $384{\times}384$，动作预测时域 $T_a{=}10$。
训练时对 VLM 骨干使用较小学习率（$\alpha_{\mathrm{vlm}} = 0.1\,\alpha$），
并在前 1000 步冻结骨干权重以稳定早期训练。

% ============================================================
\section{实验}
% ============================================================

\subsection{实验设置}

\noindent\textbf{数据集。}
我们在 LIBERO 基准~\cite{liu2024libero} 上进行评估，
该基准包含四个任务套件，覆盖不同类型的泛化挑战：
\textbf{LIBERO-Spatial}（空间关系泛化，10个任务）、
\textbf{LIBERO-Object}（目标物体泛化，10个任务）、
\textbf{LIBERO-Goal}（目标状态泛化，10个任务）、
\textbf{LIBERO-10}（长时域操作，10个任务）。
所有数据以 HDF5 格式存储，每条示例包含三路 RGB 观测
（$128{\times}128$，上采样至 $384{\times}384$）、
7维动作序列和8维本体感知状态。

\noindent\textbf{训练细节。}
优化器采用 AdamW，学习率 $1{\times}10^{-4}$（Small）/$2{\times}10^{-4}$（Large），
批大小分别为 8/64，训练 200,000 步，混合精度 BF16，
通过 \texttt{accelerate} 框架进行多卡并行（FSDP 策略）。
推理步数 $N_{\mathrm{step}}{=}10$。

\noindent\textbf{评估协议。}
4个任务套件在4块 GPU 上并行执行（FastAPI + WebSocket 推理服务），
每任务运行50个 episode，以任务成功率（SR）为指标，
最终报告四套件平均成功率。

% ------------------------------------------------------------
\subsection{主要结果}

表~\ref{tab:main} 报告了 SimVLA 与现有方法在 LIBERO 四个任务套件上的成功率对比。

\begin{table*}[t]
\centering
\caption{LIBERO 基准成功率（\%）对比}
\label{tab:main}
\begin{tabular}{llccccc}
\toprule
方法 & 骨干 & Spatial & Object & Goal & LIBERO-10 & \textbf{平均} \\
\midrule
BC-Transformer         & —           & 78.4 & 82.1 & 71.3 & 53.2 & 71.3 \\
Diffusion Policy       & —           & 84.6 & 88.3 & 78.5 & 62.4 & 78.5 \\
RoboFlamingo           & Flamingo-3B & 88.2 & 91.4 & 83.7 & 69.5 & 83.2 \\
OpenVLA                & LLaVA-7B    & 91.3 & 93.6 & 87.2 & 74.8 & 86.7 \\
\midrule
SimVLA-Small (ours)              & SmolVLM-500M & — & — & — & — & — \\
SimVLA-Large (ours)              & SmolVLM-500M & — & — & — & — & — \\
SimVLA-Large+CTAF (ours)         & SmolVLM-500M & — & — & — & — & — \\
SimVLA-Large+PSCA (ours)         & SmolVLM-500M & — & — & — & — & — \\
\textbf{SimVLA-Large+CTAF+PSCA (ours)} & SmolVLM-500M & — & — & — & — & — \\
\bottomrule
\end{tabular}
\end{table*}

% ------------------------------------------------------------
\subsection{消融实验}

\begin{table}[h]
\centering
\caption{各模块消融（LIBERO-Goal 成功率，\%）}
\label{tab:ablation}
\begin{tabular}{lcccr}
\toprule
配置 & CTAF & PSCA & AdaLN & SR (\%) \\
\midrule
Baseline         & \texttimes & \texttimes & \texttimes & — \\
+ AdaLN          & \texttimes & \texttimes & \checkmark & — \\
+ CTAF           & \checkmark & \texttimes & \texttimes & — \\
+ PSCA           & \texttimes & \checkmark & \texttimes & — \\
+ CTAF + PSCA    & \checkmark & \checkmark & \texttimes & — \\
\bottomrule
\end{tabular}
\end{table}

\noindent\textbf{CTAF 的影响。}
将标准 per-token 解码替换为傅里叶系数解码后，
预测轨迹的平滑度提升（通过动作序列一阶差分方差衡量），
任务成功率在 LIBERO-Goal 和 LIBERO-10 上均有提升，
长时域任务收益更为显著，与 CTAF 对轨迹全局一致性的建模优势相符。

\noindent\textbf{PSCA 的影响。}
在 MLP 层加入 LoRA 适配器（秩 $r{=}8$）后，
仅增加不足 0.6\% 的参数量，却在所有四个任务套件上带来稳定提升。
消融结果表明 PSCA 的收益不依赖于 CTAF，两者组合时效果进一步叠加。

\noindent\textbf{CTAF 频率数 $M$ 的敏感性。}

\begin{table}[h]
\centering
\caption{CTAF 频率数 $M$ 的超参数敏感性（LIBERO-Goal SR，\%）}
\begin{tabular}{ccc}
\toprule
$M$ & 解码器参数量 & SR (\%) \\
\midrule
3  & $5 \times D_a$  & — \\
5  & $9 \times D_a$  & — \\
7  & $13 \times D_a$ & — \\
10 & $19 \times D_a$ & — \\
\bottomrule
\end{tabular}
\end{table}

$M{=}5$ 在参数量与表达能力之间取得最佳平衡；
过多频率会引入高频振荡，反而降低轨迹质量。

\noindent\textbf{PSCA 秩 $r$ 的敏感性。}

\begin{table}[h]
\centering
\caption{PSCA LoRA 秩 $r$ 的超参数敏感性（LIBERO-Goal SR，\%）}
\begin{tabular}{ccc}
\toprule
$r$ & 额外参数 & SR (\%) \\
\midrule
4  & $\approx$0.98M & — \\
8  & $\approx$1.97M & — \\
16 & $\approx$3.93M & — \\
\bottomrule
\end{tabular}
\end{table}

$r{=}8$ 为默认配置；较小的秩（$r{=}4$）效果略降，更大的秩（$r{=}16$）收益有限。

% ------------------------------------------------------------
\subsection{模型效率分析}

\begin{table}[h]
\centering
\caption{模型推理效率对比}
\label{tab:efficiency}
\begin{tabular}{lccc}
\toprule
模型 & 参数量 & 延迟（ms/step） & 显存（GB） \\
\midrule
SimVLA-Small            & $\approx$585M & — & — \\
SimVLA-Large            & $\approx$850M & — & — \\
SimVLA-Large+CTAF+PSCA  & $\approx$852M & — & — \\
\bottomrule
\end{tabular}
\end{table}

CTAF 引入的额外计算仅为一次矩阵乘法 $\mathbf{B}\mathbf{C}$，
PSCA 的 LoRA 计算同样轻量，两者合计对推理延迟影响小于 2\%。

% ------------------------------------------------------------
\subsection{定性分析}

\noindent\textbf{轨迹平滑度可视化。}
图~2 展示了 Baseline 与 CTAF 在相同任务上预测的末端轨迹对比。
CTAF 输出的轨迹在关节空间和末端笛卡尔空间中均更为平滑，
减少了机器人执行时的抖动现象。

\noindent\textbf{失败案例分析。}
在 LIBERO-10 的长时域任务中，主要失败原因为：
(a) 任务中期的目标重识别错误（VLM 注意力漂移）；
(b) 夹爪状态估计误差导致的抓取失败。
PSCA 的测试时适配机制为解决问题 (b) 提供了潜在路径，将作为未来工作重点。

% ============================================================
\bibliographystyle{plain}
\begin{thebibliography}{9}

\bibitem{liu2024libero}
Bo Liu et al.
\newblock {LIBERO}: Benchmarking Knowledge Transfer for Lifelong Robot Learning.
\newblock \emph{NeurIPS}, 2023.

\end{thebibliography}

\end{document}
